%% PSW Dispersion Simulation — Lab Report (Notion Dark Style)
%% Compiler: xelatex
\documentclass[11pt, a4paper]{article}

%% ── 한글 + 유니코드 ────────────────────────────────────────────────────────
\usepackage{fontspec}
\usepackage{kotex}

%% ── 색상·레이아웃 ─────────────────────────────────────────────────────────
\usepackage{xcolor}
\usepackage{geometry}
\usepackage{setspace}

%% ── 수식 ─────────────────────────────────────────────────────────────────
\usepackage{amsmath}

%% ── 표 ───────────────────────────────────────────────────────────────────
\usepackage{booktabs}
\usepackage{colortbl}
\usepackage{array}
\usepackage{tabularx}

%% ── 코드 ─────────────────────────────────────────────────────────────────
\usepackage{listings}

%% ── 링크 ─────────────────────────────────────────────────────────────────
\usepackage[hidelinks]{hyperref}

%% ── 리스트 ───────────────────────────────────────────────────────────────
\usepackage{enumitem}

%% ── 다이어그램 ───────────────────────────────────────────────────────────
\usepackage{tikz}
\usetikzlibrary{shapes.geometric, arrows.meta, positioning, calc, fit}

%% ── 박스·콜아웃 ──────────────────────────────────────────────────────────
\usepackage[most, breakable]{tcolorbox}
\tcbuselibrary{skins}

%% ── 헤더/제목 스타일 ─────────────────────────────────────────────────────
\usepackage{fancyhdr}
\usepackage{titlesec}
\usepackage{parskip}

%% ═══════════════════════════════════════════════════════════════════════
%%  Notion 다크 팔레트
%% ═══════════════════════════════════════════════════════════════════════
\definecolor{nbg}{HTML}{191919}
\definecolor{ntext}{HTML}{E9E9E7}
\definecolor{nmuted}{HTML}{9B9A97}
\definecolor{nh1}{HTML}{FFFFFF}
\definecolor{ncallout}{HTML}{252525}
\definecolor{nborder}{HTML}{424242}
\definecolor{ncodebg}{HTML}{1E1E1E}
\definecolor{ncodetext}{HTML}{D4D4D4}
\definecolor{nkw}{HTML}{569CD6}
\definecolor{ncmt}{HTML}{6A9955}
\definecolor{nstr}{HTML}{CE9178}
\definecolor{ntblhd}{HTML}{2A2A2A}
\definecolor{ntblalt}{HTML}{1F1F1F}
\definecolor{naccent}{HTML}{5B89C3}
\definecolor{nlink}{HTML}{7EB3FF}
\definecolor{nhr}{HTML}{3B3B3B}
\definecolor{nyellow}{HTML}{E6A817}
\definecolor{ngreen}{HTML}{60BE8A}
\definecolor{nred}{HTML}{E06C75}
\definecolor{npurple}{HTML}{C586C0}

%% ═══════════════════════════════════════════════════════════════════════
%%  페이지 & 폰트
%% ═══════════════════════════════════════════════════════════════════════
\geometry{a4paper, top=2cm, bottom=2.5cm, left=2.8cm, right=2.8cm,
          headheight=15pt, headsep=10pt}

\pagecolor{nbg}
\color{ntext}
\setlength{\parindent}{0pt}
\setlength{\parskip}{5pt}

%% Windows 폰트 (XeLaTeX)
\setmainfont{Malgun Gothic}[
  BoldFont  = {Malgun Gothic Bold},
  Scale     = 1.0
]
\setmonofont{Consolas}[Scale=0.88]
\newfontfamily\emojifont{Segoe UI Emoji}[Scale=0.95]
\newcommand{\emo}[1]{{\emojifont #1}}

%% ═══════════════════════════════════════════════════════════════════════
%%  Hyperref
%% ═══════════════════════════════════════════════════════════════════════
\hypersetup{
  colorlinks=true,
  linkcolor=nlink, urlcolor=nlink, citecolor=nlink,
}

%% ═══════════════════════════════════════════════════════════════════════
%%  제목 스타일
%% ═══════════════════════════════════════════════════════════════════════
\titlespacing*{\section}      {0pt}{1.8em}{0.4em}
\titlespacing*{\subsection}   {0pt}{1.1em}{0.25em}
\titlespacing*{\subsubsection}{0pt}{0.7em}{0.15em}

\titleformat{\section}[block]
  {\color{nh1}\large\bfseries}
  {}
  {0em}
  {}
  [\vspace{0.1em}{\color{nhr}\hrule\vspace{0.2em}}]

\titleformat{\subsection}[block]
  {\color{ntext}\normalsize\bfseries}
  {}
  {0em}
  {}

\titleformat{\subsubsection}[block]
  {\color{nmuted}\normalsize}
  {}
  {0em}
  {}

%% ═══════════════════════════════════════════════════════════════════════
%%  콜아웃 박스
%% ═══════════════════════════════════════════════════════════════════════
\newtcolorbox{callout}[1][📌]{%
  enhanced, breakable,
  colback=ncallout, colframe=naccent,
  boxrule=0pt, leftrule=3pt,
  arc=3pt, outer arc=3pt,
  left=8pt, right=6pt, top=5pt, bottom=5pt,
  fontupper=\small\color{ntext},
  before upper={{\emojifont #1}\enspace},
}
\newtcolorbox{tip}[1][💡]{%
  enhanced, breakable,
  colback=ncallout, colframe=nyellow,
  boxrule=0pt, leftrule=3pt,
  arc=3pt,
  left=8pt, right=6pt, top=5pt, bottom=5pt,
  fontupper=\small\color{ntext},
  before upper={{\emojifont #1}\enspace},
}
\newtcolorbox{warn}[1][⚠️]{%
  enhanced, breakable,
  colback=ncallout, colframe=nred,
  boxrule=0pt, leftrule=3pt,
  arc=3pt,
  left=8pt, right=6pt, top=5pt, bottom=5pt,
  fontupper=\small\color{ntext},
  before upper={{\emojifont #1}\enspace},
}
\newtcolorbox{perf}[1][⚡]{%
  enhanced, breakable,
  colback=ncallout, colframe=ngreen,
  boxrule=0pt, leftrule=3pt,
  arc=3pt,
  left=8pt, right=6pt, top=5pt, bottom=5pt,
  fontupper=\small\color{ntext},
  before upper={{\emojifont #1}\enspace},
}

%% ═══════════════════════════════════════════════════════════════════════
%%  코드 블록 (VS Code 다크 테마)
%% ═══════════════════════════════════════════════════════════════════════
\lstdefinestyle{nd}{%
  backgroundcolor  = \color{ncodebg},
  commentstyle     = \color{ncmt}\itshape,
  keywordstyle     = \color{nkw}\bfseries,
  numberstyle      = \tiny\color{nmuted},
  stringstyle      = \color{nstr},
  basicstyle       = \ttfamily\footnotesize\color{ncodetext},
  breakatwhitespace= false,
  breaklines       = true,
  keepspaces       = true,
  numbers          = left,
  numbersep        = 8pt,
  showspaces       = false,
  showstringspaces = false,
  tabsize          = 4,
  language         = Matlab,
  frame            = single,
  rulecolor        = \color{nborder},
  xleftmargin      = 1.5em,
  framexleftmargin = 1.5em,
  aboveskip        = 0.5em,
  belowskip        = 0.3em,
  morekeywords     = {parfor, parcluster, parpool, gpuArray, gpuDevice,
                      gather, gcp, saveProfile, vertcat, sub2ind,
                      fullfile, strrep, sprintf, fprintf, repmat, ceil},
}
\lstset{style=nd}

%% ═══════════════════════════════════════════════════════════════════════
%%  표 스타일
%% ═══════════════════════════════════════════════════════════════════════
\arrayrulecolor{nborder}
\renewcommand{\arraystretch}{1.35}
\newcolumntype{L}[1]{>{\raggedright\arraybackslash}p{#1}}
\newcolumntype{C}[1]{>{\centering\arraybackslash}p{#1}}

%% ═══════════════════════════════════════════════════════════════════════
%%  헤더/푸터
%% ═══════════════════════════════════════════════════════════════════════
\pagestyle{fancy}
\fancyhf{}
\fancyhead[L]{\small\color{nmuted} PSW Dispersion Simulation}
\fancyhead[R]{\small\color{nmuted} 2026년 4월}
\fancyfoot[C]{\color{nmuted}\small\thepage}
\renewcommand{\headrulewidth}{0pt}

%% ═══════════════════════════════════════════════════════════════════════
%%  TikZ 플로우차트
%% ═══════════════════════════════════════════════════════════════════════
\tikzset{
  fnode/.style={%
    rectangle, rounded corners=4pt,
    draw=nborder, fill=ncallout,
    text=ntext,
    minimum width=7cm, minimum height=0.68cm,
    align=center, font=\small,
    inner sep=5pt,
  },
  farrow/.style={-Stealth, color=nmuted, line width=1.2pt},
}

%% ════════════════════════════════════════════════════════════════════════
\begin{document}
%% ════════════════════════════════════════════════════════════════════════

%% ── 제목 ──────────────────────────────────────────────────────────────
\vspace*{-0.3em}
\begin{center}
  {\Huge\bfseries\color{nh1}%
   \emo{🧲}\enspace PSW Dispersion Simulation}\\[0.5em]
  {\large\color{nmuted} 병렬화 최적화 · 기술 참조 문서}\\[0.15em]
  {\normalsize\color{nmuted} 2026년 4월 · i7-13700K / RTX 4060 Ti / DDR5 64 GB}
\end{center}
\vspace{0.5em}{\color{nhr}\hrule}\vspace{1em}

%% ── 개요 콜아웃 ────────────────────────────────────────────────────────
\begin{callout}[🧲]
PSW(Polarized Spin-Wave) Dispersion Simulation 전체 문서 허브.
\texttt{run\_all.m}을 기준으로 각 모듈별 설명으로 구성된다.
이미지 처리(image processing) 관련 내용은 \textbf{PSW 이미지 처리 분석} 별도 문서로 분리되어 있다.
\end{callout}

\vspace{1.2em}

%% ── 문서 구성 ─────────────────────────────────────────────────────────
{\color{nh1}\large\bfseries 문서 구성}
\vspace{0.5em}

\begin{center}
\begin{tabular}{>{\color{naccent}\small}l >{\small\color{ntext}}L{9.5cm}}
\rowcolor{ntblhd}
\multicolumn{1}{l}{\color{nmuted}\small 페이지} &
\multicolumn{1}{l}{\color{nmuted}\small 내용} \\
\hline
1.\ 개요 및 설정 & 병렬 전략, Parallel Pool, 상수/파라미터 설정 \\
\rowcolor{ntblalt}
2.\ 시뮬레이션 핵심 로직 & linearly.mat 로드, 레이저 세기 계산, parfor sub-batch 루프 \\
3.\ 후처리 및 결과 저장 & 결과 재조립, .mat 저장 \\
\rowcolor{ntblalt}
2-1.\ time\_integrate.m 분석 & Velocity Verlet 적분, GPU/CPU 경로, axay\_opt\_partial 호출 구조 \\
\end{tabular}
\end{center}

\vspace{1.5em}

%% ── 전체 실행 흐름 플로우차트 ─────────────────────────────────────────
{\color{nh1}\large\bfseries 전체 실행 흐름}
\vspace{0.8em}

\begin{center}
\begin{tikzpicture}[node distance=0.55cm]
  \node[fnode] (A) {\texttt{run\_all.m} 시작};
  \node[fnode, below=of A] (B) {Parallel Pool 설정 (8 workers)};
  \node[fnode, below=of B] (C) {STN\_constants + linearly.mat 로드};
  \node[fnode, below=of C] (D) {레이저 세기 $I_1$, $I_2$ 계산};
  \node[fnode, below=of D] (E) {parfor: 80 태스크 병렬 실행};
  \node[fnode, below=of E] (F) {\texttt{time\_integrate.m} — Velocity Verlet CPU 적분};
  \node[fnode, below=of F] (G) {결과 재조립 \texttt{Cell\_velocity}};
  \node[fnode, below=of G] (H) {\texttt{result\_*.mat} 저장};
  \foreach \f/\t in {A/B, B/C, C/D, D/E, E/F, F/G, G/H}
    \draw[farrow] (\f) -- (\t);
\end{tikzpicture}
\end{center}

\newpage

%% ════════════════════════════════════════════════════════════════════════
\section{\emo{⚙️}\enspace 1. 개요 및 설정}

\begin{callout}[📌]
\texttt{run\_all.m}의 Section 0--3을 다룬다: Parallel Pool 설정, STN 상수 로드,
레이저 세기 파라미터 계산.
\end{callout}

\subsection{Section 0 — Parallel Pool 설정}

MATLAB local profile의 \texttt{MaxNumWorkers} 기본값은 물리 P-core 수(8)로 제한된다.
\texttt{parcluster} API를 통해 직접 프로파일을 수정해야 원하는 worker 수를 사용할 수 있다.

\begin{lstlisting}
n_workers = 8;   % 실측 최적값

p = gcp('nocreate');
if ~isempty(p) && p.NumWorkers ~= n_workers
    delete(p);  p = [];
end
if isempty(p)
    cluster = parcluster('local');
    if cluster.NumWorkers ~= n_workers
        cluster.NumWorkers = n_workers;
        saveProfile(cluster);   % 프로파일에 영구 저장
    end
    parpool(cluster, n_workers);
    fprintf('Parallel pool 시작: %d workers\n', n_workers);
end
\end{lstlisting}

\subsection{Section 1 — STN 상수 로드 및 MAX 오버라이드}

\begin{lstlisting}
load(fullfile(root_dir, 'STN_constants.mat'));
MAX = 2000000;   % STN_constants의 MAX=20,000 -> 200,000 오버라이드
\end{lstlisting}

\begin{tip}[💡]
\texttt{STN\_constants.mat}에는 \texttt{MAX=20000}이 저장되어 있으나
시뮬레이션은 200,000 입자가 필요하므로 로드 직후 오버라이드한다.
\end{tip}

\subsection{Section 3 — 양자상태 파일 경로}

\begin{lstlisting}
jm_data_path = fullfile(root_dir, 'data', jm_data);
if ~isfile(jm_data_path)
    error('파일 없음: %s', jm_data_path);
end
\end{lstlisting}

%% ════════════════════════════════════════════════════════════════════════
\section{\emo{🔬}\enspace 2. 시뮬레이션 핵심 로직}

\begin{callout}[📌]
\texttt{run\_all.m}의 Section 4--6을 다룬다:
linearly.mat 로드, 레이저 세기 계산, parfor sub-batch 루프.
\end{callout}

\subsection{Section 4 — linearly.mat 로드 및 2D reshape}

\begin{lstlisting}
load(fullfile(root_dir, 'linearly.mat'));
[nj_lin, nm_lin, ~] = size(linearly);
linearly_2d = reshape(linearly, nj_lin*nm_lin, []);
% (nj*nm) x nI -- 모든 parfor worker에 broadcast
clear linearly;
\end{lstlisting}

\begin{tip}[💡]
\texttt{linearly.mat}는 $\langle\cos^2\theta\rangle$ 룩업 테이블이다.
\texttt{reshape}으로 2D로 펼친 뒤 broadcast 변수로 사전 로드하면,
각 worker가 매 task마다 파일을 다시 읽는 오버헤드를 없앤다.
\end{tip}

\subsection{Section 5 — 레이저 세기 사전 계산}

\begin{lstlisting}
max_I  = 2.0 / (pi * w0z * w0y) * sqrt(4.0 * log(2) / (pi * tau^2));
max_I  = max_I * atten1 * atten2 * 1e-3;
I1_val  = max_I * I1_list;       % 스칼라
I2_vals = max_I * I2_list(:);    % n_groups x 1
\end{lstlisting}

\subsection{Section 6 — parfor Sub-Batch 병렬 루프}

\subsubsection{Sub-batch 분해 전략}

\[
  \underbrace{200{,}000}_{\text{MAX/조건}} \times \underbrace{8}_{\text{n\_groups}}
  \;\longrightarrow\;
  \underbrace{20{,}000}_{\text{batch\_size}} \times \underbrace{80}_{\text{n\_tasks}}
\]

각 태스크 번호 \texttt{task} $\in [1,\,80]$에서 조건 인덱스 매핑:
\[
  g = \left\lceil \frac{\texttt{task}}{n_\text{sub}} \right\rceil
\]

\begin{lstlisting}
batch_size = 20000;
n_sub      = ceil(MAX / batch_size);   % 10
n_tasks    = n_groups * n_sub;          % 80
results    = cell(n_tasks, 1);

parfor task = 1:n_tasks
    g_idx = ceil(task / n_sub);

    [particle_t, laser_t] = init_particles(jm_data_path, batch_size, ...
        x_position_fwhm, Dye_sigma, Dye_tau, ...
        vx_caused_tilted, x_init_vel_sigma, y_init_vel_sigma, ...
        z_init_vel_off, z_init_vel_sigma, Ij_distribution, Ips_z, Ips_y);

    j_idx_t = particle_t(:,10) + 1;
    m_idx_t = particle_t(:,11) + 1;
    aligned_cost_t = linearly_2d(sub2ind([nj_lin,nm_lin], j_idx_t,m_idx_t), :);

    I1_t = repmat(I1_val,         batch_size, 1);
    I2_t = repmat(I2_vals(g_idx), batch_size, 1);

    results{task} = time_integrate( ...
        particle_t(:,1:3), particle_t(:,4:6), ...
        laser_t(:,3), laser_t(:,4:5), ...
        I1_t, I2_t, aligned_cost_t, zeros(batch_size,2), ...
        lambda, w0z, w0y, tau, batch_size, tstep, t_ini, tend, false);
end
\end{lstlisting}

%% ════════════════════════════════════════════════════════════════════════
\section{\emo{📊}\enspace 3. 후처리 및 결과 저장}

\begin{callout}[📌]
\texttt{run\_all.m}의 후처리 부분(Section 7, 9)을 다룬다.
sub-batch 결과 재조립 → 결과 저장 순서로 진행된다.
이미지 처리(Section 8)는 \textbf{PSW 이미지 처리 분석} 문서로 분리됨.
\end{callout}

\subsection{Section 7 — Sub-batch 결과 재조립}

80개 태스크의 결과를 조건($I_2$ 세기)별로 합쳐 \texttt{Cell\_velocity}에 저장한다.

\begin{lstlisting}
Cell_velocity = cell(n_groups, 1);
for g = 1:n_groups
    task_idx = (g-1)*n_sub + (1:n_sub);
    Cell_velocity{g} = vertcat(results{task_idx});  % MAX x 3
end
\end{lstlisting}

\begin{tip}[💡]
\texttt{vertcat(results\{task\_idx\})}는 셀 배열 슬라이싱과 수직 결합을 한 번에 처리한다.
각 cell은 $\texttt{batch\_size}\times 3$ 행렬이므로 결합 후 $\texttt{MAX}\times 3$ 행렬이 된다.
\end{tip}

\subsection{Section 9 — 결과 저장}

결과 파일명은 \texttt{jm\_data} 파일명 기반으로 자동 생성된다.
저장 전에 \texttt{constants} 구조체를 먼저 구성한다.

\subsubsection{constants 구조체 구성}

\begin{lstlisting}
% 시뮬레이션 제어
constants.n_workers     = n_workers;
constants.MAX           = MAX;
constants.batch_size    = batch_size;
% 이미지 처리 파라미터 (-> PSW 이미지 처리 분석 문서 참조)
constants.rep_std       = rep_std;
constants.rep_vel       = rep_vel;
constants.vel_pix       = vel_pix;
constants.atten_factorx = atten_factorx;
constants.atten_factory = atten_factory;
% 분자빔 초기조건
constants.jm_data           = jm_data;
constants.x_position_fwhm   = x_position_fwhm;
constants.Dye_sigma          = Dye_sigma;
constants.Dye_tau            = Dye_tau;
constants.vx_caused_tilted   = vx_caused_tilted;
constants.x_init_vel_sigma   = x_init_vel_sigma;
constants.y_init_vel_sigma   = y_init_vel_sigma;
constants.z_init_vel_off     = z_init_vel_off;
constants.z_init_vel_sigma   = z_init_vel_sigma;
constants.Ij_distribution    = Ij_distribution;
constants.Ips_z              = Ips_z;
constants.Ips_y              = Ips_y;
% 레이저 초기조건
constants.lambda = lambda;
constants.w0z    = w0z;
constants.w0y    = w0y;
constants.tau    = tau;
constants.atten1 = atten1;
constants.atten2 = atten2;
constants.tstep  = tstep;
constants.t_ini  = t_ini;
constants.tend   = tend;
\end{lstlisting}

\subsubsection{저장 명령}

\begin{lstlisting}
out_file = fullfile(root_dir, ...
    sprintf('result_%s.mat', strrep(jm_data, '.txt', '')));
save(out_file, 'Cell_velocity', 'I1_list', 'I2_list', 'constants');
fprintf('저장 완료: %s\n', out_file);
\end{lstlisting}

\subsubsection{저장 변수 목록}

\vspace{0.3em}
\begin{center}
\begin{tabular}{>{\ttfamily\small\color{naccent}}l >{\small\color{ntext}}L{9cm}}
\rowcolor{ntblhd}
\multicolumn{1}{l}{\color{nmuted}\small 변수} &
\multicolumn{1}{l}{\color{nmuted}\small 설명} \\
\hline
Cell\_velocity  & 조건별 최종 속도 데이터 (n\_groups$\times$1 cell, 각 MAX$\times$3) \\
\rowcolor{ntblalt}
I1\_list, I2\_list & 사용된 레이저 세기 목록 \\
constants & 시뮬레이션 전체 파라미터 구조체 — 제어(n\_workers, MAX, batch\_size), 분자빔 초기조건, 레이저 초기조건, 이미지 처리 파라미터 포함 \\
\end{tabular}
\end{center}

\subsection{성능 벤치마크}

\vspace{0.3em}
\begin{center}
\begin{tabular}{>{\small}L{5.5cm} >{\small\centering}C{2.5cm} >{\small\centering\arraybackslash}C{2.5cm}}
\rowcolor{ntblhd}
\color{nmuted}\small 방식 &
\color{nmuted}\small 소요 시간 &
\color{nmuted}\small 배속 \\
\hline
순차 실행 & $\sim$960 s & 1$\times$ \\
\rowcolor{ntblalt}
\bfseries\color{ngreen} parfor 8 workers & \bfseries\color{ngreen} $\sim$181 s & \bfseries\color{ngreen} 5.3$\times$ \\
parfor 12 workers & $\sim$183 s & 5.2$\times$ \\
\rowcolor{ntblalt}
parfor 16 workers & $\sim$187 s & 5.1$\times$ \\
parfor 24 workers & $\sim$208 s & 4.6$\times$ \\
\end{tabular}
\end{center}
\vspace{0.4em}

\begin{perf}[⚡]
8 workers가 최적값이다. Worker 수를 늘릴수록 \texttt{aligned\_cost} 배열(broadcast 변수)의
동시 접근이 DDR5 메모리 버스를 포화시켜 오히려 성능이 저하된다.
\end{perf}

\newpage

%% ════════════════════════════════════════════════════════════════════════
\section{\emo{⚡}\enspace 2-1. time\_integrate.m 분석}

\begin{callout}[📌]
\texttt{time\_integrate.m}은 Velocity Verlet 시간 적분기로, GPU/CPU 두 경로를 모두 지원한다.
\end{callout}

\subsection{함수 시그니처}

\begin{lstlisting}
function velocity = time_integrate(position, velocity, Ij, Ips, ...
        I1_vec, I2_vec, aligned_cost, last_data, ...
        lambda, w0z, w0y, tau, MAX, tstep, t_ini, tend, use_gpu)
% use_gpu:  true  -> GPU 사용 (단독 실행 시)
%           false -> CPU 사용 (parfor worker 내 권장)
%           생략  -> GPU 자동 감지
\end{lstlisting}

\subsection{CPU 경로 (parfor worker 내부)}

\begin{lstlisting}
t = t_ini;
while t < tend
    position = position + velocity .* tstep;
    [~, last_data, a, ~] = axay_opt_partial( ...
        position, t + Ij, I1_vec, I2_vec, ...
        aligned_cost, last_data, Ips, lambda, w0z, w0y, tau, MAX);
    velocity = velocity + a .* tstep;
    t = t + tstep;
end
\end{lstlisting}

\begin{tip}[💡]
총 6,000 스텝 ($t_\text{ini}=-30\,\text{ns}$, $t_\text{end}=+30\,\text{ns}$, $\Delta t=10\,\text{ps}$).
각 스텝에서 위치 업데이트 → 광학 힘 계산 → 속도 업데이트 순서로 진행된다.
\end{tip}

\subsection{GPU 경로 (단독 실행 시)}

\begin{lstlisting}
g_dev = gpuDevice;
nI = size(aligned_cost, 2);
bytes_per_particle = nI * 8 + 250;
gpu_batch_size = floor(g_dev.AvailableMemory * 0.8 / bytes_per_particle);
gpu_batch_size = min(gpu_batch_size, MAX);
% VRAM 가용 메모리 기반 자동 청킹
\end{lstlisting}

\begin{warn}[⚠️]
현재 parfor worker 내부에서는 \texttt{use\_gpu=false}로 CPU 모드 강제 사용.
20,000 입자 규모에서는 GPU kernel launch overhead가 이득을 상쇄한다
(실측: GPU 84~s vs CPU 64~s).
\end{warn}

%% ════════════════════════════════════════════════════════════════════════
\section{\emo{🔧}\enspace 병렬화 전략 상세}

\subsection{MaxNumWorkers 강제 확장}

MATLAB local profile은 \texttt{MaxNumWorkers} 기본값이 8이다.
단순히 \texttt{parpool(N)}만 호출하면 8개 이상 연결되지 않는다.

\begin{lstlisting}
cluster = parcluster('local');
if cluster.NumWorkers ~= n_workers
    cluster.NumWorkers = n_workers;
    saveProfile(cluster);   % 프로파일에 영구 저장
end
parpool(cluster, n_workers);
\end{lstlisting}

\subsection{핵심 최적화 요약}

\vspace{0.3em}
\begin{center}
\begin{tabular}{>{\small\bfseries\color{naccent}}L{4cm} >{\small\color{ntext}}L{9cm}}
\rowcolor{ntblhd}
\multicolumn{1}{l}{\color{nmuted}\small 최적화 항목} &
\multicolumn{1}{l}{\color{nmuted}\small 내용} \\
\hline
aligned\_cost 벡터화 &
  \texttt{for} 루프(MAX회) $\to$ \texttt{reshape}+\texttt{sub2ind} 단일 연산으로 대체 \\
\rowcolor{ntblalt}
gauss\_t 중복 제거 &
  시간 루프 내 6회 중복 계산 $\exp(-4\ln 2\cdot t^2/\tau^2) \to$ 1회 사전 계산 \\
broadcast 변수 &
  \texttt{linearly\_2d}를 parfor 외부에서 한 번 로드, 모든 worker에 공유 \\
\rowcolor{ntblalt}
col\_idx 경계 처리 &
  \texttt{min(I\_floor+1, n\_cols-1)}로 룩업 테이블 out-of-range 방어 \\
GPU 호환 할당 &
  \texttt{zeros(MAX,3,'like',position)} — CPU/GPU 경로 공통 코드 유지 \\
\end{tabular}
\end{center}

\subsection{알려진 한계 및 개선 방안}

\vspace{0.3em}
\begin{center}
\begin{tabular}{>{\small\color{nmuted}}L{3.5cm} >{\small\color{ntext}}L{9.5cm}}
\rowcolor{ntblhd}
\multicolumn{1}{l}{\color{nmuted}\small 항목} &
\multicolumn{1}{l}{\color{nmuted}\small 내용} \\
\hline
RAM 사용률 & $\sim$85\% (64 GB 기준 $\sim$54 GB): batch\_size 증가 불가 \\
\rowcolor{ntblalt}
메모리 대역폭 & 8 workers가 포화점 → 추가 병렬화 불가 \\
Single precision 전환 &
  \texttt{linearly\_2d} + \texttt{aligned\_cost}를 \texttt{single}로 변환 시
  메모리 85\% $\to$ 55\%, 속도 10--15\% 향상 예상 \\
\rowcolor{ntblalt}
GPU 재검토 & 입자 수 $>100{,}000$/batch 규모에서 재평가 가능 \\
\end{tabular}
\end{center}

\vspace{1em}
{\color{nhr}\hrule}
\vspace{0.5em}
{\small\color{nmuted}
  이미지 처리(전자 반동 노이즈, 2D Voigt PSF, hist3 합성곱, 프로젝션)는
  \textbf{PSW 이미지 처리 분석} 문서를 참조.
}

\end{document}
